% =====================================================================
%  CHƯƠNG 9. SỰ GIAO THOA GIỮA LLM AGENT VÀ REINFORCEMENT LEARNING
% =====================================================================
\chapter{Sự giao thoa giữa LLM Agent và học tăng cường}\label{ch:giaothoa}

\section{Học tăng cường từ phản hồi con người (RLHF)}
Các tác tử AI hiện đại dựa trên LLM liên hệ chặt chẽ với học tăng cường qua kỹ thuật \textbf{RLHF} (Reinforcement Learning
from Human Feedback) \cite{christiano_rlhf}. Quy trình gồm ba bước (Hình~\ref{fig:ch9-rlhf}):
\begin{enumerate}
  \item \textbf{Huấn luyện mô hình phần thưởng (Reward Model).} Thu thập đánh giá ưu tiên của con người giữa các câu trả
        lời của mô hình (câu nào tốt hơn), từ đó huấn luyện một hàm phần thưởng $R_\phi(x, y)$ chấm điểm câu trả lời $y$
        cho câu hỏi $x$.
  \item \textbf{Tối ưu chính sách bằng PPO.} Dùng thuật toán Actor-Critic PPO (Proximal Policy Optimization) cập nhật
        trọng số LLM để tối đa hóa điểm của Reward Model, đồng thời phạt độ lệch KL so với mô hình gốc để mô hình không
        ``lách luật'' phần thưởng:
        \[
          \max_\theta\; \mathbb{E}_{x,\, y \sim \pi_\theta}\!\left[ R_\phi(x,y) \right]
          - \beta\, \mathrm{KL}\!\left( \pi_\theta(\cdot \mid x) \,\|\, \pi_{\text{ref}}(\cdot \mid x) \right).
        \]
  \item \textbf{Định hình hành vi an toàn.} RLHF biến một mô hình dự đoán từ kế tiếp thô thành một tác tử hội thoại biết
        từ chối yêu cầu độc hại, tuân thủ chỉ dẫn và giải trình mạch lạc.
\end{enumerate}
Có thể đặt RLHF vào khung MDP của Chương~4: trạng thái là câu hỏi cùng phần câu trả lời đã sinh, hành động là từ kế tiếp,
phần thưởng do Reward Model chấm khi câu trả lời hoàn tất. Đây là cùng một vòng lặp tác tử -- môi trường mà tác tử
Q-Learning của bài thực hành sử dụng, chỉ khác ở quy mô.

\hinhtikz{ch9_rlhf}{Quy trình ba bước của RLHF: dữ liệu ưu tiên của con người, mô hình phần thưởng và tối ưu chính sách
bằng PPO có ràng buộc KL.}{ch9-rlhf}

\section{Tác tử tự hoàn thiện và suy luận theo cây MCTS}
Các hệ tác tử tiên tiến đưa thuật toán học tăng cường kinh điển vào chính quá trình suy luận (inference-time compute):
\begin{itemize}
  \item \textbf{Monte Carlo Tree Search (MCTS) và Q* Search.} Thay vì sinh một luồng suy luận duy nhất, tác tử khám phá
        một cây suy luận nhiều nhánh qua bốn bước chọn -- mở rộng -- đánh giá -- lan truyền ngược. Một mô hình đánh giá
        từng bước (Process Reward Model -- PRM) chấm điểm mỗi bước tư duy, đóng vai hàm giá trị giống như $Q(s,a)$
        trong Q-Learning (Hình~\ref{fig:ch9-mcts}).
  \item \textbf{Tự sửa sai (Self-Correction).} Tác tử tự tạo môi trường phản hồi (self-play): viết mã, chạy kiểm thử đơn vị,
        đọc lỗi và sửa, lặp lại cho tới khi đạt mục tiêu. Kết quả kiểm thử chính là tín hiệu phần thưởng.
\end{itemize}

\hinhtikz{ch9_mcts}{Suy luận dạng cây MCTS với mô hình đánh giá từng bước (PRM).}{ch9-mcts}

\section{Xu thế hội tụ Neuro-Symbolic trong hệ đa tác tử}
Kết hợp năng lực lập luận ngôn ngữ của LLM (Neuro) với các giải thuật tối ưu hóa điều khiển chính xác của RL và quy hoạch
động (Symbolic/Dynamic Programming) mở ra các hệ thống tự hành mới:
\begin{itemize}
  \item \textbf{Phân tầng điều khiển (Hierarchical Agent Control).} Tác tử LLM cấp cao làm Meta-Planner: phân tích mục tiêu
        nghiệp vụ, chia thành mục tiêu con và giao cho các tác tử RL chuyên biệt thực thi thao tác vi mô thời gian thực
        (robot, giao dịch tần suất cao) (Hình~\ref{fig:ch9-hier}).
  \item \textbf{Trao đổi thông điệp chuẩn hóa.} Các tác tử trong Agentic Mesh tương tác qua giao thức mở, bảo đảm tính
        minh bạch và khả năng kiểm toán toàn diện như yêu cầu quản trị ở Chương~\ref{ch:quantri}.
\end{itemize}

\hinhtikz{ch9_hierarchical}{Phân tầng điều khiển: tác tử LLM Meta-Planner giao mục tiêu con cho các tác tử RL chuyên
biệt.}{ch9-hier}
